% Demo paper for the meertex website screenshots, as it was submitted. Every
% name, number and reference in it is made up. revision/main.tex is the same
% paper after the reviews.
\documentclass[acmsmall,review,anonymous]{acmart}

\settopmatter{printacmref=false}
\setcopyright{none}
\renewcommand\footnotetextcopyrightpermission[1]{}
\acmJournal{TOSEM}

\begin{document}

\title{Reproducible Energy Measurements for Software Engineering Research}

\author{Lena Vogt}
\affiliation{\institution{Example University}\country{Germany}}
\email{lena@example.com}

\begin{abstract}
Energy measurements of software are hard to repeat. Results change with the
machine, its temperature, the background load and the measurement tool. We
reviewed 120 studies that report the energy use of software and found that
few of them report enough detail to repeat the measurement. We propose a
checklist of twelve items and show on three case studies that it makes
measurements repeatable within 3\,\%.
\end{abstract}

\maketitle

\section{Introduction}

Software engineering research increasingly reports the energy use of
programs, libraries and development practices. A result that cannot be
repeated on another machine says little about the software and a lot about
the setup it was measured on.

\section{Literature Review}

We searched four digital libraries for studies that measure the energy of
software and kept 120 of them. For each study we recorded the hardware, the
measurement tool, the number of repetitions and whether idle power was
subtracted. Few studies report all four.

\section{The Checklist}

The checklist has twelve items in three groups. The first group describes
the machine, the second the measurement and the third the analysis. Each item
is a question a reader must be able to answer from the paper alone.

\section{Case Studies}

We applied the checklist to three earlier studies and repeated their
measurements on a second machine. With the checklist the results agreed
within 3\,\%. Without it they differed by up to 40\,\%.

\section{Conclusion}

Repeatable energy measurements need a description of the setup that is as
careful as the description of the software under test. The checklist is a
first step towards it.

\end{document}
